% Page 043 — page imprimée 39
% Contenu seul : le préambule, l'en-tête et la pagination sont gérés par meyrueis.tex

\sectiondate{Bonnes pages des souvenirs de Jean Gounelle (1902-1994)}{1933-1935}

\begin{center}
\vspace{0.4em}
\textbf{Pasteur de Meyrueis de 1932 à 1935}
\end{center}

Avec Meyrueis le jeune ménage pastoral que nous formions, Irène et moi, abordait un
pays plus civilisé que le Pompidou. Le territoire paroissial qui, désormais, allait nous
incomber comportait une petite ville et trois vallées.

Meyrueis constituait alors, comme aujourd’hui encore, un centre avec docteur,
dentiste, PTT, pharmacien, garage, Caisse d’Epargne, Banque, et beaucoup de
magasins etc. pourvu en outre de nombreux hôtels, dont certains luxueux. On sentait
la petite cité touristique, soignée, attirante, se suffisant à elle même, susceptible
d’accueillir une clientèle exigeante.

La population de ce bourg se rattachait, en grande partie, à l’Eglise Catholique. Les
protestants, majoritaires au 19\textsuperscript{e} siècle, ayant les premiers et assez massivement quitté
la ville, les Caussenards venaient combler les vides ainsi créés. L’école laïque ne
conservait qu’une influence secondaire par rapport à l’école libre, tenue par de jeunes
frères, instituteurs en soutane, très dynamiques et animés d’un grand zèle. De plus s’y
trouvait implantée une maison de « sœurs des pauvres » Ces femmes, d’un
dévouement sans borne, travaillaient dans ce milieu, appréciées de tous.
Notre ministère où la partie médicale avait tenu un si grand rôle au Pompidou, ne
pouvait qu’en être profondément modifié. Ce côté cependant, n’en fut pas pour
autant négligeable. !!

Notre action s’engagea surtout dans les trois vallées, entièrement protestantes elles,
celle de la Jonte, de la Brèze et du Bétuzon, trois ruisseaux dont les noms poétiques
chantaient comme le bruissement de leurs eaux. Ces trois vallées se rejoignaient à
l’endroit même où n’avait pu que surgir, dans la nuit des temps, un groupement
humain, protégé, cohérent et naturel.
Je ne vois pas d’autre explication à la naissance et à l’origine de Meyrueis.

Pour rester dans les considérations générales et m’en tenir à ce qui constituait ma
charge en ce lieu, disons que la paroisse protestante, dans les années 1930, avait un
double visage. Celui de l’hiver dans l’axe de ces trois vallées où se trouvaient nos
paysans, beaucoup plus à l’aise que ceux du Pompidou, à cause de leur lait de brebis
exploité par les caves de Roquefort. Celui de l’été. De Juin à Octobre, l’auditoire au
culte, nombreux et divers, se composait surtout de fidèles issus de grandes paroisses
citadines. Le Château d’Ayres, en particulier, dont le propriétaire faisait partie du
Conseil Presbytéral, constituait un pôle d’attraction exceptionnel.

Nos familles bourgeoises protestantes de Paris et d’ailleurs venaient volontiers y
passer l’été et revenaient ensuite. Elles animaient, nos « ventes », nos fêtes d’Eglise
organisaient au Temple de nombreuses rencontres et parfois de magnifiques concerts.
Le pasteur devait donc s’adapter à ce double aspect de l’Eglise locale, et ne pas
négliger les uns aux dépens des autres. Il convient également de se souvenir qu’en ces temps-là le tourisme constituait un luxe réservé à quelques privilégiés et que les
vacances étaient surtout le fait d’une élite et non pas généralisées comme aujourd’hui.
Les problèmes avaient donc peu de rapport avec ceux rencontrés au Pompidou, sauf
qu’il fallait, à travers un complexe différent en ces deux centres lozériens, faire
entendre le même message de paix, d’amour, de lumière et de justice: le message de
l’Evangile.[\ldots]

\subsection{Juillet 1932 : Installation à Meyrueis}

Ce mois de juillet fut, cette année là, particulièrement pluvieux. C’est sous les
parapluies et dans la boue que s’effectuèrent nos premiers contacts avec la
population. Notre adaptation ne fut certainement pas facilitée par la clémence du
temps ! J’avais peine à croire que le dépassement de la ligne de partage des eaux qui
séparait le Pompidou de Meyrueis pût amener de tels contrastes de climat. La
Cévenne elle même différait profondément. L’absence de châtaigniers ici était totale
(sauf dans les deux vallées schisteuses et granitiques du Bétuzon et de la Brèze, voire au dessus du
château d’Ayres). Tourné vers l’Atlantique, le paysage se montrait humide, plus
verdoyant, moins abrupt avec cependant, des entailles profondes à la séparation des
Causses, préfigurant déjà les Gorges du Tarn. Bref un changement plus accentué que
prévu.

Cette première impression plutôt fâcheuse fut compensée par notre prise de
possession d’un presbytère vraiment confortable pour l’époque.
Mon prédécesseur le pasteur Albert Chaudier, l’avait laissé en parfait état. Il avait
même voulu, en guise de bienvenue, que les pièces fussent retapissées à neuf pour
nous. Cela nous toucha beaucoup. Quel progrès sur ce que nous quittions ! Eau
courante, cheminées dans toutes les pièces au rez-de-chaussée et au premier,
harmonie dans la disposition de l’ensemble. La cuisine, la salle de séjour et le bureau
donnaient par plusieurs sorties sur un agréable jardin, hautement clôturé, que
remplissaient deux immenses poiriers chargés de fruits. Un conseiller presbytéral
facétieux, me disait « Il conviendrait d’appeler votre demeure « la villa des poires ».
Le moment n’était pas encore venu de posséder une salle de bains installée ni de
toilettes intérieures. Mais ces dernières se trouvaient à proximité et tout pouvait se
passer dans la plus normale discrétion. !

Très conscients de cette situation améliorée, quelle fut notre première préoccupation ?
Aussi inattendue que soit la réponse, la vérité m’oblige à la donner : mettre en route
un enfant.

[ La grossesse fut délicate et le premier-né vit le jour à la Maison de Santé de Nîmes. On l’appela
André, il devint pasteur... comme son père, puis professeur à la faculté de Théologie de Montpellier
écrivit de nombreux livres... ]

Mon ministère fut renouvelé par le bonheur de cette naissance... Je visitai d’arrache-
pied. J’organisai aussi parfaitement que possible mon enseignement auprès des
enfants et des jeunes, écoles du dimanche, du jeudi, catéchismes, Union Chrétienne.
Je montai une chorale et animai des réunions le soir. Je suppléai ainsi à l’absence
momentanée de collaboration de ma femme, absorbée, au début tout au moins, par
son bébé. Mais elle reprit assez vite sa place. Nous avions pu nous faire aider par la
fille d’un conseiller presbytéral, Jeanne Poujol (fille des Poujol, fermiers de la Fabrique) qui se mit à notre service et nous épaula avec beaucoup de dévouement. Les jours de
veillées ou de réunions, Irène pouvait m’accompagner. Au retour, dans la nuit, je
ramenai Jeanne chez elle.

Il y eut, au cours de notre séjour, ici et là, quelques soins médicaux à donner. Bien
plus tard, venant d’Afrique du Nord en vacances et de passage dans la région, tel ou
tel ancien paroissien de Meyrueis ou de Gatuzières nous disait : « Vous souvenez-
vous, Madame Gounelle vous avez soigné ma fille », « vous avez guéri la grand
mère », » « vous avez avant notre départ donné le landau d’André ce qui nous a
soulagés » « Quand ma mère, alitée depuis longtemps, souffrait de nombreuses
escarres, vous avez arrangé son lit et amélioré sa position en sorte que ses
souffrances furent considérablement allégées mais une heure plus tard elle
recommença à souffrir et me disais : Si notre dame était là ! »
Ce côté de notre mission restait bien réel. Mais au Pompidou, évidemment, sans
docteur et sans bonnes sœurs la trace s’était montrée bien plus profonde.
C’est ainsi que, sans éclat particulier, dans la fidélité aux petites choses, se déroula un
ministère paroissial normal, sans heurts et sans secousses, pendant trois ans.

\subsection{Notre concierge Monsieur Avesque (dit Lou Rascle)}

Un autre trait particulier du Meyrueis de 1930 c’est la personnalité de notre
concierge. Il s’appelait monsieur Avesque. Grand et original, fin parfois dans ses
réflexions, obstiné dans son comportement, vulgaire trop souvent dans son langage, il
avait le don d’imposer sa manière de faire, sans nuance. Très âgé, il avait dû, dans sa
jeunesse être un vrai colosse. Courbé par les ans, il restait encore très vigoureux.
Chaque semaine, il passait une journée entière, pour une rémunération dérisoire, à
nettoyer ce temple immense, briquer ses pavés disjoints, enlever la poussière de ses
bancs tous numérotés, garnir l’hiver les trois grands poêles, bien incapables d’ailleurs
de fournir une chaleur suffisante. Aussi chacun portait son chauffe-pied pour le
service du dimanche, qu’il avait préparé aussi bien que possible.

Après avoir allumé les feux ( il disait éclairé le feu ), le moment venu, il s’attaquait à
la lourde cloche, au son grave et profond, qui se répercutait au loin dans les vallées,
trois fois dans la matinée. Le dernier coup indiquait le début de l’office. Très souvent,
l’église catholique annonçait sa messe à la même heure. M. Avesque n’admettait pas
que nos « frères séparés » (ainsi les désignait-il) mettent leur cloche en branle tant
qu’il n’avait pas terminé. S’il avait commencé le premier, on devait respecter sa
priorité. Sinon il n’arrêtait pas de tirer sur sa corde. Cela aurait pu durer une demi-
heure, ou davantage, il ne cédait pas souvent, impuissant devant cette guerre des
cloches, j’ai été contraint de retarder « ma cérémonie » (sic) car il fallait bien, que le
silence se fit.

\subsection{Vie de l’église}

L’école du dimanche précédant le culte réunissait de nombreux élèves, pas toujours
très disciplinés. Je tenais beaucoup à ce qu’il y règne une certaine liberté et que cette
école ne ressemble pas, cet égard, à l’école de chaque jour. Nous avions plusieurs
groupes avant l’assemblée générale. Après le rassemblement, quand M. Avesque
s’apercevait qu’un garçon chahutait un peu trop à son gré, ne me trouvant pas assez
sévère, il faisait chauffer son tisonnier à blanc, puis, en pleine leçon il l’approchait, menaçant, de la figure du garnement, au risque de provoquer un accident. Les enfants
se tenaient tranquilles ; terrifiés par le concierge. Plusieurs, devenus vieux
aujourd’hui, en gardent le souvenir ! Ils me l’ont confirmé il y a peu de temps !
\emph{(je revois encore, en 1940, le Rascle nous poursuivant dans le temple avec son tisonnier, car nous le
faisions enrager ; chauffé à blanc ? bien sûr que non ! se souvient Paul Loupiac)}

Quant aux baptêmes, mariages et enterrements, il exigeait que les familles, après
s’être entretenues avec le pasteur, aillent le trouver, et surtout ne l’oublient pas sous
peine de ne pouvoir accéder au temple, qu’il estimait être un peu sa propriété
personnelle. En cas de négligence, il menaçait de prendre le large à l’heure de la
cérémonie, et de se rendre avec les clefs dans sa poche dans son champ le plus
éloigné. Cela arriva au moins une fois et n’amusa personne !!!

Irène tenait habituellement l’harmonium. Mais, avant et après la naissance d’André,
elle dut renoncer provisoirement à cette charge. Elle fut remplacée par une jeune fille
pleine de bonne volonté, mais qui ne pouvait accompagner les cantiques que d’une
main. L’instrument était bon, un harmonium de quatre jeux et demi, tout à fait adapté
à l’édifice. Un dimanche où je mettais au point avec un groupe de fidèles et
l’organiste, la partie musicale, Mr. Avesque s’approcha de nous et dit :
« \emph{Mademoiselle Agulhon, vous êtes beaucoup plus forte que le pasteur. Lui il a
besoin de deux mains pour jouer, vous, vous vous en tirez avec un doigt !} » On
pouvait considérer qu’il s’agissait d’une réflexion amusante dont on aurait dû
seulement sourire, mais cela déplut visiblement à l’organiste !

A cette période vivait à Meyrueis un pasteur à la retraite, originaire du pays, tout à
fait digne et remarquable. Il se nommait David Poujol\footnote{David Poujol était l’oncle de Marcel Poujol dit « Poujol d’Alger », rue saint blaise, décédé en 1994.}. Le concierge avait usé avec
lui ses fonds de culotte sur les bancs de l'école communale de Meyrueis. Ce pasteur
avait dirigé des églises importantes, Mazamet, Nîmes, La Haye et d’autres lieux.
Orateur incontestable, doué d’une puissante voix, on le sentait habitué aux nombreux
rassemblements. Il cultivait l’éloquence classique à laquelle il excellait. Un
dimanche, alors qu’il montait en chaire (où il me remplaçait, accompagné du
concierge qui le tutoyait depuis toujours, ce dernier lui dit devant moi : « \emph{Alors
David, c’est toi qui vas faire l’article aujourd’hui ?} » ce qui lui attira une verte
réplique, méritée, qui lui cloua le bec !
Je désire rendre hommage à ce pasteur distingué, disparu depuis longtemps, et qui
imposait le respect. Il finissait ses jours à Meyrueis où il vivait dans sa modeste
maison veuf et seul. Au gros de l’hiver il se rendait chez une de ses filles mais
revenait très vite. C’était un ramasseur de champignons. [\ldots]

Puisque je parle de David Poujol, comment ne pas mentionner le rayonnement de
deux autres pasteurs qui passaient, à cette époque leurs vacances d’été à Meyrueis.
Emile Guiraud de Paris et Louis Hubac de Castres. Ce dernier était le fils d’un paysan
que nous appréciions beaucoup, habitant la vallée de la Brèze, propriétaire de la
ferme de Conillergues près de Campis. Lorsqu’il prêchait, le temple se remplissait.
Sa prédication originale ( je me souviens d’un sermon sur « la main ») attirait, mais
n’était-il pas, surtout, un enfant du pays ?

Il passait régulièrement au presbytère, une ou deux fois par semaine, pour me
demander si j’étais libre. Il m’est arrivé assez souvent de l’accompagner. J’ai ainsi
pris contact, grâce à lui, avec une population catholique assez arriérée, de l’avis
même du docteur, mais qui s’est toujours montrée heureuse et même honorée d’être
visitée par un pasteur, dépourvu, de toute évidence, d’esprit de propagande ou de
prosélytisme.

A propos de ces traversées des Causses j’ai présenté à la mémoire un voyage en car
que j’y fis un jour, je ne sais plus pour quelle raison. Au départ de Meyrueis, au
dernier moment, un jeune curé vint s’asseoir à côté de moi, très exubérant et gentil
au possible, il engagea avec moi une conversation intéressante et animée. Au fur et à
mesure de l’entretien, une véritable sympathie s’établit entre nous deux. Quelle
agréable compagnie pour ce déplacement ! Puis, m’interrogeant plus directement, il
voulut savoir qui j’étais. Je n’avais aucune raison de le lui cacher : le pasteur de
Meyrueis. L’effet fut immédiat. Il me tourna carrément le dos et plus un mot ne sortit
de sa bouche. Quand il descendit , il ne me dit même pas au revoir !
Je n’ai jamais oublié cette attitude blessante, et il m’est arrivé, plus tard au cours de
réunions œcuméniques au Maroc, par exemple, avec d’autres curés, de raconter ce
petit incident pénible pour me réjouir qu’aujourd’hui, ce genre de comportement ne
serait ni possible ni concevable.

Même si des désaccords doctrinaux subsistent, au moins nous avons vis à vis les uns
des autres une attitude fraternelle grâce au Concile Vatican II et au pape Jean XXIII.

Enfin un mot sur les concerts. Il y en eut beaucoup au cours de notre séjour. Les plus
appréciés furent de deux sortes :

Ceux organisés avec le concours du célèbre violoncelliste de l’époque, d’origine
cévenole Jacques Serres. Il vient de mourir au moment où j’écris ces lignes. J’avais
béni, à Valleraugue son mariage avec la pianiste Ady Levastré. C’était vraiment du
grand art et le milieu estivant, très mélomane se trouvait à même de l’apprécier.
J’ajoute que le milieu paysan ne se montrait pas insensible au talent de cet artiste.

Je l’accompagnai à l’harmonium avec le souci évident de le mettre en valeur, mais en
m’efforçant quand même de dégager toutes les ressources de cet instrument, trop
décrié, de nos temples ou églises de montagne. Ensuite, de Meyrueis, nous étendîmes
notre audience à d’autres paroisses de la région, puis jusqu’à Pau, Millau, Montauban
et plus tard l’Afrique du nord

La seconde série de concerts se fit, en particulier, avec le concours d’une personne
qui faisait partie des vacanciers habituels, la Comtesse de Trémeuge. Cette femme
jeune, assez extraordinaire d’une grande beauté, avait reçu de la nature une voix
admirable. Elle n’appartenait pas à notre église, mais aimait fréquenter notre
presbytère. Elle admirait le dépouillement de la foi réformée. Devant cette artiste,
séduit par sa valeur et sa personne, je lui avais conseillé de lui permettre au moins du Bach du
Haendel ou du Franck. Je lui avais fait répéter, à la maison, des extraits de la Passion
selon St-Matthieu. Elle s’arrêta net et modestement, j’ajoute sagement, elle me
conseilla de lui permettre de chanter nos beaux cantiques, soigneusement
sélectionnés. Le succès fut total.

Quant au pasteur Guiraud, incomparable orateur, doué d’une voix plus distinguée que
forte, quel enchantement que d’entendre cet homme ! Je me souviens d’une prière de
lui, entendue à Meyrueis, qui commençait par cette invocation : « \emph{O Toi qui donnes,
Dieu de Noël, O Toi qui pardonnes, Dieu de Pâques, nous venons à Toi !} [\ldots]

Pendant deux ans, les circonstances locales s’y prêtant, je fis la classe de 6\textsuperscript{e} et de 5\textsuperscript{e} à
une jeune fille très douée, dont les parents, excellents paroissiens, ne voulaient pas se
séparer. Elle s’appelait Jeanine Fabre.
Je me fis donc professeur. Ma femme et moi prîmes en charge toutes les matières du
programme : français, latin, anglais, histoire, géographie, botanique géologie etc. La
suite de ces leçons elle fut admise avec félicitations dans un lycée de la capitale. Outre
la joie que l’on éprouve quand on a à faire à une enfant intelligente, cela apporta un
peu d’aisance à notre foyer qui en avait bien besoin. La contre-partie, un peu fâcheuse
c’est que toutes nos matinées furent prises. Autre conséquence : un bruit favorable se
répandit autour de ces leçons, et des vacanciers voulurent également que, au cours de
l’été, je fasse travailler leurs enfants. Je ne m’y suis pas refusé.

Janine Fabre, après y avoir fait ses études, fut nommée professeur de lettres au Lycée
Victor Duruy dans le 7\textsuperscript{e} Arrondissement.

\begin{itemize}
\item \emph{Plus tard elle épousa Francis Nougarède, ingénieur géologue. Ils eurent trois enfants,
Florence, Olivier et Pascal qui revenaient ensemble ou à tour de rôle au « Clair Logis »
leur maison de Meyrueis qui avait auparavant appartenu à la tante de Janine, le Dr. Juliette
Cabanel,
Francis fut emporté par une tumeur en 2004 et Janine le suivit un an plus tard.}
\end{itemize}

\emph{« Le Clair Logis » est passé désormais en d’autres mains.}

Nous avions aussi au mois de juillet et d’août étudiants en Théologie ou en Lettres
comme pensionnaires. Au nombre de 3 ou 4 chaque fois. Ils nous étaient envoyés par
l’Association Libérale de Strasbourg. Nous les faisions coucher dans deux vastes
chambres au premier. Toute la journée ils couraient dans les montagnes, visitaient la
région et, le soir venu, ils entonnaient des cantiques au presbytère. Cela tenait lieu de
culte de famille. Quel ne fut pas notre étonnement de constater que des
rassemblements assez importants de touristes se formaient autour de notre demeure
pour les écouter chaque jour ! Ces voix d’hommes étaient très belles et le sens musical
de ces jeunes strasbourgeois incontestable !

Comment ne pas mentionner nos rapports avec le Docteur Lapeyre et mes escapades
en sa compagnie sur les Causses. Dans une petite cité les rapports entre le pasteur et
Le médecin sont très importants, à tous égards. A Meyrueis, le docteur, catholique
pratiquant, préférait la compagnie du pasteur à celle du prêtre. Il l’avait déjà montré
avec un de mes prédécesseurs, M. Maurin.
Sa clientèle, en dehors du centre et des vallées, s’étendait sur les Causses, territoire
dont les habitants m’étaient inconnus, car je n’y avais aucun paroissien. Son rayon
d’action atteignait la Malène et Sainte Enimie, sur le Tarn, juste avant les Gorges
proprement dites. n’aimait pas rouler seul dans sa voiture, une belle conduite
intérieure Citroën confortable.


\subsection{Climat général en 1932 -- 1935}

De profonds changements se dessinèrent au cours de ces trois années. Les mutations
ne datent pas d’aujourd’hui seulement... la radio s’étendait partout et prenait une
grande place dans la vie de nos campagnes. Nous avions acquis un petit poste « Point
Bleu » bien insuffisant et modeste, mais nous l’écoutions tous les jours. Grâce à une
antenne tendue depuis le clocher du temple jusqu’au toit du presbytère, en l’absence
de moteurs et d’appareils ménagers, je captais les émissions avec une grande pureté.
C’était l’époque où les vedettes se nommaient Gaby Morlay, Sim Viva ou la
danseuse Mitsi, beautés lointaines et inconnues mais derrière le nom desquelles
j’imaginais la grâce féminine parisienne. L’époque également, hélas, des grands
scandales financiers, Stavisky et autres.

Puis ce fut le 6 Février 1934 et surtout le début des hurlements sinistres d’un inconnu
qui se nommait Adolf Hitler. Vraiment j’avais le sentiment que la relative tranquillité
que nous avions connue depuis l’armistice de 1919 était, bel et bien, en train de
disparaître.

Voilà un éclairage bien incomplet de ce temps de Meyrueis. Je pourrais développer.
Parler de Millau, de Florac, de nos veillées à Gatuzières ou à Cabrillac, de nos
campagnes « Christianisme Social », de nos représentations théâtrales avec les jeunes
etc Tout cela fut fort intéressant, remarquable même. Mais pourquoi le temps passé à
Meyrueis fut-il si court ?...

\emph{[... ( Jean Gounelle, comme plusieurs membres de sa famille, était attiré par l’Afrique du Nord :
l’énigme de ces pays évangélisés aux premiers siècles, leur essor, puis l’élimination à peu près totale
des églises chrétiennes par l’Islam du 8\textsuperscript{o} au 12\textsuperscript{o} siècle... Ne serait-ce pas le moment de s’y
réimplanter ? Oran n’avait plus de pasteur... Il fut nommé... et partit sans attendre... )....]}

Mais, cinquante ans plus tard, je conserve au fond de mon cœur, empreint d’une
certaine nostalgie, l’apport précieux que fut pour nous notre bref séjour à Meyrueis.
Et quand j’y retourne maintenant de Mandagout, (\emph{près du Vigan}) je perçois mieux, me
semble-t-il, les richesses et les bénédictions reçues dans cette humble cité ! Je m’en
souviens comme on se souvient avec émotion de ces périodes de préparation intense,
aussi bien intérieure que sur le terrain, où, loin des apparences et guidé par une main
invisible, se joue vraiment le destin d’un homme.
